% Generated draft: requires mathematical and source review.
\documentclass[11pt]{article}
\usepackage[margin=25mm]{geometry}
\usepackage{fontspec}
\setmainfont{Latin Modern Roman}
\usepackage{amsmath,amssymb,mathtools,unicode-math}
\setmathfont{Latin Modern Math}
\usepackage{enumitem,needspace,xurl,hyperref}
\setlength{\parindent}{0pt}
\setlength{\parskip}{5pt}
\setlength{\emergencystretch}{3em}
\newcommand{\N}{\mathbb{N}}
\newcommand{\Z}{\mathbb{Z}}
\newcommand{\Q}{\mathbb{Q}}
\newcommand{\R}{\mathbb{R}}
\newcommand{\C}{\mathbb{C}}
\newcommand{\F}{\mathbb{F}}
\newcommand{\A}{\mathbb{A}}
\newcommand{\PP}{\mathbb{P}}
\newcommand{\E}{\mathbb{E}}
\newcommand{\eps}{\varepsilon}
\newcommand{\dd}{\,\mathrm{d}}
\begin{document}
\section*{UnsolvedMath ID: 801099}
\textbf{Source problem code:} catalogue entry
\subsection{Definitions and mathematical statement}
The supplied record refers to a ring \(\mathcal{O}_{K,S}\), an ideal sequence \((\mathfrak{a}_n)_{n\geq 1}\), a parameter \(d\), and an ideal norm \(\operatorname{Norm}\), but does not define these objects or their underlying dynamical construction. In particular, the hypotheses on \(K\), \(S\), and \(d\), and the construction of \(\mathfrak{a}_n\), remain unresolved. The following therefore preserves the source questions as a conditional formulation, rather than supplying a fully specified problem.

Determine the circumstances under which
\[
\log\operatorname{Norm}(\mathfrak{a}_n)=o(d^n)
\qquad (n\to\infty).
\]
Whenever the logarithm and quotient are defined, this asks for conditions ensuring
\[
\lim_{n\to\infty}\frac{\log\operatorname{Norm}(\mathfrak{a}_n)}{d^n}=0.
\]
Also determine the circumstances under which there exists an ideal \(\mathfrak{b}\subset\mathcal{O}_{K,S}\) such that
\[
\#\{n\in\Z_{\geq 1}:\mathfrak{a}_n\mid\mathfrak{b}\}=\infty.
\]
The source additionally states that
\[
\mathfrak{a}_1\mid\mathfrak{a}_n\qquad\text{for every }n\geq 1,
\]
and identifies \(\mathfrak{b}=\mathfrak{a}_1\) as a reasonable candidate. This candidate is a suggestion in the source, not an asserted conclusion. Under the usual convention for integral ideals, \(I\mid J\) means \(J=IC\) for some integral ideal \(C\); the intended ideal class and norm convention must be confirmed from the missing context.
\subsection{Short English statement}
Characterize when the source's ideal sequence has logarithmic norm growing more slowly than \(d^n\), and when infinitely many terms divide one fixed ideal of \(\mathcal{O}_{K,S}\). The source suggests the first term as that fixed ideal. The sequence, norm, parameter, and ambient hypotheses are missing, so the intended H\'enon-map problem cannot yet be stated completely.
\subsection{Sources}
\begin{itemize}
\item Local UnsolvedMath registry, permanent ID \(801099\), exact source code \texttt{catalogue entry}, title ``Henon Maps: Open Problems --- Question 51.'' The supplied registry attributes the problem to Berger and Bedford, with proposed year 2023, and labels it open; this status is supplied metadata, not an independently verified assessment.
\item Berger and Bedford, \emph{Henon Maps: A List of Open Problems} (2023), Question 51. Supplied source URL: \url{https://arxiv.org/abs/2312.03907}. The record reports access on 2026-10-16 and extraction by \texttt{arxiv\_latex}; these are supplied provenance details, not a claim of independent source access or verification.
\end{itemize}
\end{document}
